% Options for packages loaded elsewhere
% Options for packages loaded elsewhere
\PassOptionsToPackage{unicode}{hyperref}
\PassOptionsToPackage{hyphens}{url}
\PassOptionsToPackage{dvipsnames,svgnames,x11names}{xcolor}
%
\documentclass[
  12pt,
]{article}
\usepackage{xcolor}
\usepackage[margin=1in]{geometry}
\usepackage{amsmath,amssymb}
\setcounter{secnumdepth}{5}
\usepackage{iftex}
\ifPDFTeX
  \usepackage[T1]{fontenc}
  \usepackage[utf8]{inputenc}
  \usepackage{textcomp} % provide euro and other symbols
\else % if luatex or xetex
  \usepackage{unicode-math} % this also loads fontspec
  \defaultfontfeatures{Scale=MatchLowercase}
  \defaultfontfeatures[\rmfamily]{Ligatures=TeX,Scale=1}
\fi
\usepackage{lmodern}
\ifPDFTeX\else
  % xetex/luatex font selection
\fi
% Use upquote if available, for straight quotes in verbatim environments
\IfFileExists{upquote.sty}{\usepackage{upquote}}{}
\IfFileExists{microtype.sty}{% use microtype if available
  \usepackage[]{microtype}
  \UseMicrotypeSet[protrusion]{basicmath} % disable protrusion for tt fonts
}{}
\makeatletter
\@ifundefined{KOMAClassName}{% if non-KOMA class
  \IfFileExists{parskip.sty}{%
    \usepackage{parskip}
  }{% else
    \setlength{\parindent}{0pt}
    \setlength{\parskip}{6pt plus 2pt minus 1pt}}
}{% if KOMA class
  \KOMAoptions{parskip=half}}
\makeatother
% Make \paragraph and \subparagraph free-standing
\makeatletter
\ifx\paragraph\undefined\else
  \let\oldparagraph\paragraph
  \renewcommand{\paragraph}{
    \@ifstar
      \xxxParagraphStar
      \xxxParagraphNoStar
  }
  \newcommand{\xxxParagraphStar}[1]{\oldparagraph*{#1}\mbox{}}
  \newcommand{\xxxParagraphNoStar}[1]{\oldparagraph{#1}\mbox{}}
\fi
\ifx\subparagraph\undefined\else
  \let\oldsubparagraph\subparagraph
  \renewcommand{\subparagraph}{
    \@ifstar
      \xxxSubParagraphStar
      \xxxSubParagraphNoStar
  }
  \newcommand{\xxxSubParagraphStar}[1]{\oldsubparagraph*{#1}\mbox{}}
  \newcommand{\xxxSubParagraphNoStar}[1]{\oldsubparagraph{#1}\mbox{}}
\fi
\makeatother


\usepackage{longtable,booktabs,array}
\usepackage{calc} % for calculating minipage widths
% Correct order of tables after \paragraph or \subparagraph
\usepackage{etoolbox}
\makeatletter
\patchcmd\longtable{\par}{\if@noskipsec\mbox{}\fi\par}{}{}
\makeatother
% Allow footnotes in longtable head/foot
\IfFileExists{footnotehyper.sty}{\usepackage{footnotehyper}}{\usepackage{footnote}}
\makesavenoteenv{longtable}
\usepackage{graphicx}
\makeatletter
\newsavebox\pandoc@box
\newcommand*\pandocbounded[1]{% scales image to fit in text height/width
  \sbox\pandoc@box{#1}%
  \Gscale@div\@tempa{\textheight}{\dimexpr\ht\pandoc@box+\dp\pandoc@box\relax}%
  \Gscale@div\@tempb{\linewidth}{\wd\pandoc@box}%
  \ifdim\@tempb\p@<\@tempa\p@\let\@tempa\@tempb\fi% select the smaller of both
  \ifdim\@tempa\p@<\p@\scalebox{\@tempa}{\usebox\pandoc@box}%
  \else\usebox{\pandoc@box}%
  \fi%
}
% Set default figure placement to htbp
\def\fps@figure{htbp}
\makeatother





\setlength{\emergencystretch}{3em} % prevent overfull lines

\providecommand{\tightlist}{%
  \setlength{\itemsep}{0pt}\setlength{\parskip}{0pt}}



 
\usepackage[]{biblatex}


\usepackage{amsthm}
\usepackage{amsmath}
\usepackage{amsfonts}
\usepackage{amssymb}
\usepackage{float}
\usepackage{caption}
\usepackage{subcaption}
\usepackage{stmaryrd}
% Numbering is tied to sections (e.g., Proposition 2.1, 2.2, etc.)
% Theorems and Propositions (italic)
\theoremstyle{plain}
\newtheorem{theorem}{Theorem}[section]
\newtheorem{proposition}[theorem]{Proposition}
\newtheorem{lemma}[theorem]{Lemma}
% Definitions and Corollaries (upright)
\theoremstyle{definition}
\newtheorem{definition}{Definition}
\newtheorem{corollary}{Corollary}
\newtheorem{example}{Example}
% \newtheorem{algorithm}[theorem]{Algorithm}
% Optional: formatting for the proof environment
\renewenvironment{proof}
   {\par\noindent\textbf{Proof.}\ }
   {\hfill$\blacksquare$\par}
% \usepackage{algorithm}
% \usepackage{algorithmic}
\usepackage[ruled,vlined,linesnumbered]{algorithm2e}
\makeatletter
\@ifpackageloaded{caption}{}{\usepackage{caption}}
\AtBeginDocument{%
\ifdefined\contentsname
  \renewcommand*\contentsname{Table of contents}
\else
  \newcommand\contentsname{Table of contents}
\fi
\ifdefined\listfigurename
  \renewcommand*\listfigurename{List of Figures}
\else
  \newcommand\listfigurename{List of Figures}
\fi
\ifdefined\listtablename
  \renewcommand*\listtablename{List of Tables}
\else
  \newcommand\listtablename{List of Tables}
\fi
\ifdefined\figurename
  \renewcommand*\figurename{Figure}
\else
  \newcommand\figurename{Figure}
\fi
\ifdefined\tablename
  \renewcommand*\tablename{Table}
\else
  \newcommand\tablename{Table}
\fi
}
\@ifpackageloaded{float}{}{\usepackage{float}}
\floatstyle{ruled}
\@ifundefined{c@chapter}{\newfloat{codelisting}{h}{lop}}{\newfloat{codelisting}{h}{lop}[chapter]}
\floatname{codelisting}{Listing}
\newcommand*\listoflistings{\listof{codelisting}{List of Listings}}
\usepackage{amsthm}
\theoremstyle{remark}
\AtBeginDocument{\renewcommand*{\proofname}{Proof}}
\newtheorem*{remark}{Remark}
\newtheorem*{solution}{Solution}
\newtheorem{refremark}{Remark}[section]
\newtheorem{refsolution}{Solution}[section]
\makeatother
\makeatletter
\makeatother
\makeatletter
\@ifpackageloaded{caption}{}{\usepackage{caption}}
\@ifpackageloaded{subcaption}{}{\usepackage{subcaption}}
\makeatother
\usepackage{bookmark}
\IfFileExists{xurl.sty}{\usepackage{xurl}}{} % add URL line breaks if available
\urlstyle{same}
\hypersetup{
  pdftitle={Elastic Bounds of Pearson's r: Theoretical and Empirical Insights for Ordinal Variables},
  pdfauthor={Cees van der Eijk; Scott Moser},
  colorlinks=true,
  linkcolor={black},
  filecolor={Maroon},
  citecolor={RoyalBlue},
  urlcolor={BrickRed},
  pdfcreator={LaTeX via pandoc}}


\title{Elastic Bounds of Pearson's \(r\): Theoretical and Empirical
Insights for Ordinal Variables}
\author{Cees van der Eijk \and Scott Moser}
\date{28, October 2025}
\begin{document}
\maketitle

\renewcommand*\contentsname{Table of contents}
{
\hypersetup{linkcolor=}
\setcounter{tocdepth}{3}
\tableofcontents
}

\newpage

\begin{abstract}
The use of product-moment correlations to assess relationships between ordered-categorical variables is widespread and generally tolerated in spite of the non-metric character of such variables. Potential problems of this practice have been identified in the literature. This article focusses on one of these, namely that correlation coefficients are poorly comparable across pairs of ordered-categorical variables because their ranges are constrained in non-uniform ways. This not only affects the validity of correlation coefficients for descriptive purposes, but also as bases for subsequent analyses, as in, e.g., principal component analysis, factor analysis, and structural equation modelling. This article focusses on this problem of ‘constrained’ correlations. It explains the problem and explores conditions under which it is most severe. It provides illustrations based on both fabricated and actual empirical data, the latter from a mass-survey that ubiquitously contains many ordered-categorical variables that are commonly analysed with product-moment correlations. The article is accompanied by a software tool in R for analysing the attainable upper- and lower bounds of correlations in actual data. 

Pearson’s correlation coefficient $r$, when applied to ordinal data, is subject to elasticity due to constraints from marginal distributions. This paper develops theoretical bounds for $r$, based on Boole–Fréchet–Hoeffding inequalities and rearrangement inequalities, and derives analytic expressions for their limits. We investigate when symmetry in the bounds holds or breaks, and explore practical implications for correlation-based methods like PCA and SEM. An accompanying R package allows users to compute bounds and perform hypothesis tests given fixed marginals.
\end{abstract}

\newpage

\newpage

\section{Introduction and Conceptual
Setting}\label{introduction-and-conceptual-setting}

In this article we elaborate why correlations involving
ordered-categorical variables are generally `constrained', which means
that their attainable maximum and minimum values are rarely +1 and -1,
as is often thought, but values that have to be calculated for each pair
of variables. We assess the magnitude of these constraints in fictitious
as well as empirical data and analyse under which conditions they are
most severe. The article is accompanied by a software tool in R: that
allows applied researchers to easily assess the extent of the problem in
their own data.

The use of product-moment correlations (also known as Pearson
correlations) to assess linear relationships between ordered-categorical
variables --henceforth referred to as orcats-- is widespread and widely
condoned in empirical studies, in spite of such variables not being
metric while the product-moment correlation assumes that they are. We do
not focus on this problem from an obsessive desire for methodological
purity, were it only because many statistical procedures are quite
robust for violations of their assumptions {[}fn: examples?{]}. Instead,
we draw attention to problematic consequences of this violation of
assumptions that, as far as we know, have been rarely recognised in the
applied literature although they are well-known in the theoretical
statistical literature. This is the problem that correlation
coefficients involving orcats are constrained in their range which makes
them mutually incomparable, leading to problems in hypothesis testing
and in analytical procedures that use correlations as their input.

Despite long-standing cautions about applying Pearson’s correlation coefficient \(r\) to ordinal variables (\autocite{Labovitz1970}; \autocite{Gaito1980}), it remains routine in survey-based social science. In political science and sociology alike, researchers frequently correlate 4–, 5–, or 7–point items as if they were interval-scaled, usually without examining the implications of the marginals. Examples include correlations among attitude items in large-scale survey batteries (\autocite{AlwinKrosnick1991}), ideology and issue positions on mixed-resolution scales (\autocite{BrewerGross2005}; \autocite{KinderKam2009}), and summary indices that are subsequently modeled using correlation-based techniques. Reviews document that such practices are pervasive in applied survey analysis (\autocite{KreuterValliant2007}). Methodological discussions sometimes defend pragmatic use of parametric tools with ordinal data (\autocite{Jacoby1999}), but even optimistic views do not address the \emph{feasible range} of \(r\) imposed by fixed marginals—an omission with direct consequences for substantive interpretation.

Alternatives to Pearson’s \(r\)—notably Spearman’s \(\rho\), Kendall’s \(\tau\), and the polychoric correlation—offer different trade-offs (\autocite{Spearman1904}; \autocite{Kendall1938}; \autocite{Olsson1979}). Rank-based measures ignore the marginals; polychorics assume latent bivariate normality. In contrast, we pursue bounds on \(r\) that \emph{explicitly} depend on the observed marginals, drawing on rearrangement inequalities and the Fréchet–Hoeffding framework (\autocite{Frechet1951}; \autocite{Hoeffding1940}; \autocite{Rueschendorf1981}). This perspective clarifies when the attainable range of \(r\) is narrow, when symmetry fails (\(r_{\min}\neq -r_{\max}\)), and when the lower bound can be strictly positive. It also explains why correlation-based models with ordinal indicators (EFA/CFA/SEM) can be distorted when Pearson correlations are used (\autocite{FloraCurran2004}; \autocite{GreenEtAl1997}). Our contribution situates empirical practice within these theoretical limits and provides tools to compute and visualize the bounds given fixed marginals.

This article proceeds as follows. We start by briefly elaborating the
character of ordinal categorical variables (orcats), why they are
generally not regarded as being metric, and why, in spite of that, they
are nevertheless very frequently used in statistical procedures that
assume metric level data. We then proceed in section 3 to show why
correlations based on orcats tend to be constrained, which implies that
their range is smaller than from -1 to +1, and why this is problematic
in applied research. This raises the question how to determine the
actual range of attainable values of correlations between orcats, which
we address in the subsequent section. In section 5 we assess the
magnitude of these constraining effects on two different sets of data.
One is a typical survey dataset in which many of the variables are
ordinal categorical, the orcatorcat 2019 British Election Study. The
second consists of data derived from simulations covering all possible
relevant conditions for pairs of ordered categorical variables. We use
both sets of data to explore conditions that influence the magnitude of
the constraining effect. Section Section~\ref{sec-MCsim} concludes by
summarizing the findings, by elaborating the consequences for applied
research, and by discussing directions for further research.

\subsection{On ordinal-categorical variables (orcats), their level of
measurement and their common use in quantitative
analyses}\label{on-ordinal-categorical-variables-orcats-their-level-of-measurement-and-their-common-use-in-quantitative-analyses}

Orcats are ordinal variables with a relatively small number of
categories (or ranks) compared to the number of cases, thus resulting in
partial orderings of cases. Although they can be found in all field of
quantitative inquiry they are ubiquitous in data from standardised
social science surveys. Likert items and rating scales {[}describe in fn
Likert items and rating scales, add that there are yet other formats
such as semantic-differential scales that also yield orcats{]} are among
the most favoured ways to elicit responses in such surveys. These
survey-question formats have most often between 3 and 11 ordinal
categories. The numerical values assigned to the categories of orcats
are most often consecutive integers, starting with either 0 or 1. {[}fn:
explain that these values may be included in the formulation of a survey
question, which is immaterial to our discussion here. The values
assigned to the categories (and to the responses) reflect, in
combination with the analysis procedure the analyst's interpretation of
their meaning.{]} . Such an assignment of values is as good as, and
indeed analytically equivalent to, any other set of values that reflects
the analyst's understanding of the ordinal relations between the
categories, i.e., any monotonic transformation. However, objections are
often raised when orcats {[}fn: add references{]} are analysed with
procedures requiring metric measurement. That requirement implies that
the differences between values can validly be compared so that, e.g.,
the substantive difference between categories 1 and 2 is equal to that
between categories 3 and 4, and half as large as the difference between
categories 2 and 4, etc. Any implied interval-level interpretation of
ordinal categorical variables requires therefore a compelling
theoretical or empirical justification of the numerical values assigned
to the categories rather than just the analyst's fiat. In principle this
is quite possible, for example by calibrating values of categories
against psychophysical measurements (cf., Lodge; Tillie; others), or
against the results of psychometric/linguistic analysis of verbal
response labels (cf Sönning 2024). However, such procedures are rarely
integrated in applied survey research because of their high cost. In the
absence of compelling justifications for values assigned to categories
orcats can only be considered to provide ordinal information, not metric
information. In spite of this we find that they are very commonly
analysed with statistical tools that require metric measurement of
variables such as product-moment correlations or procedures that build
upon such correlations. We find this clearly exhibited in published
articles in high-ranking academic journals {[}fn: explain how we did
this{]} and in highly regarded textbooks where examples are given of the
use of correlations between orcats as input for procedures such as PCA,
FA or SEM {[}fn: examples: Kline; Byrne; xxx{]}. This practice is so
established that it is occasionally rationalised by referring to orcats
as being `quasi-interval level' measures {[}fn: reference{]}. One can,
however, also find more explicit justifications for this practice, which
emphasise the robustness of correlation coefficients for violations of
their measurement level assumptions. Labovitz (1967, 1970) in particular
has compared correlations between orcats scored according to
equal-distances with correlations when the variables were scored with a
monotonic random procedure (which only assumes ordinality between
categories). He reports that the differences are negligible, and
suggests that even if the theoretical objection against unfounded
interval level interpretations of orcats would be justified, that
objection would nevertheless have such minor practical consequences as
to make it irrelevant in most situations. Not surprisingly, Labovitz'
work has in turn be criticised itself on various grounds, most
particularly of (inadvertently) mainly covering situations where this
robustness holds, and overlooking other situations (cf O'Brien 1979;
xxx). At this place we will not focus on discussions about the
equal-distance scoring practice. Our main reason for referring to them
is to demonstrate the existence of lively debates in the applied
literature about potential issues when using orcats in procedures
requiring metric data. The existence of these debates makes it
remarkable that another important issue has --to the best of our
knowledge-- remained mainly overlooked in the applied literature while
it has been well-known in the theoretical statistical literature. This
is the problem that correlations between orcats are generally
`constrained', which means that they do not range between -1 and +1, but
between minima and maxima that have to be calculated separately for each
pair of variables. The nature of this issue is the subject of the next
section.

\subsection{Bounds of correlations between ordinal categorical variables
(orcats)}\label{bounds-of-correlations-between-ordinal-categorical-variables-orcats}

The values that product-moment correlations between ordinal-categorical
variables (orcats) can attain are generally constrained by minimum and
maximum bounds of attainable values that are not -1 and +1. {[}fn: The
maximum attainable value is +1 only when the univariate distributions of
the variables are identical. But that does not imply the minimum
attainable value then to be -1, as that also requires further conditions
that the univariate distributions ,must meet (see below, in the
discussion of Table 2). The minimum attainable value of a correlation is
-1 when the univariate distributions are each other's inverse, but
again, that does not imply that the maximum is +1, as that requires
additional conditions of the univariate distributions.{]}

Because the bounds of correlations between orcats are not -1 and +1 it
is even more difficult than commonly thought in applied research to
interpret the strength of correlations. The well-known and often
referred to criteria formulated by Cohen (1988: 77-81) for correlations
in the behavioural sciences {[}fn: Cohen suggested, with appropriate
caution, the following guidelines: \(0.10<r<0.29\) weak; \(0.30<r<0.49\)
moderate; \(0.50<r<1.00\) strong. Hemphill (2003) suggests on the basis
of a review of several meta-analyses in psychological research that the
cut-off of 0.50 for considering a correlation to be `strong' may be too
high in many common instances.{]} are explicitly based on the notion
that the minimum and maximum attainable values for r are -1 and +1. But
that is generally not the case and, moreover, the actual upper- and
lower-bounds of r are not the same across pairs of variables, but
specific to each. This state of affairs thus generates two
interconnected problems when correlating orcats. First, the absence of
clarity of the lower- and upper-bounds of r makes qualitative
interpretations of the strength of correlations elusive. Second, because
the bounds of attainable values of r are specific for each pair of
variables each correlation is expressed in terms of a different, but
unknown reference interval. Such correlations are therefore not
comparable. To better grasp these problems, it is necessary to determine
the upper- and lower bounds of correlations, given the univariate
distributions of the variables involved.

Statistical scholarship has since long recognized that, in contrast to
what is ofen thought, the upper- and lower-bounds of correlations
between orcats are generally not the reference interval {[}-1, +1{]}.
{[}insert references, including Shi\&Huang 1992; xxxx{]}. Determining
the actual bounds given the marginal distributions of the two variables
can be done by using Whitt's (1976) rearrangement theorem and
Hardy--Littlewood--Pólya's (1952) rearrangement inequality. For
obtaining the maximum attainable value of a correlation the values of X
and Y have each to be sorted independently in either ascending or
descending order and then paired (comonotonic arrangement). For the
minimum attainable value of the correlation (given the marginals of the
two variables) the anti-comonotonic arrangement is required (sorting one
of the variables in ascending and the other in descending order). A
formalized elaboration is presented in Appendix 1. Moreover, the
algorithm for calculating rmin and rmax has been programmed in R and is
{[}or: will be{]} freely available for researchers to assess the minimum
and maximum attainable values of r for pairs of variables in their own
data. Details about this routine are described in Appendix 2.

A series of illustrative cases is presented in Table 2 (as in Table 1
for n=12 and two orcats with 4 categories each). Although the
distributions of the variables given in these examples is of no
particular importance, and could have been replaced by innumerable other
examples, they are nevertheless helpful to arrive at some general
observations about rmin and rmax.:

\begin{itemize}
\tightlist
\item
  The values of rmin and rmax are not constant across all examples, and
  not even across correlations involving a specific variable (see
  examples (a) and (b) which involve the same \(X\)-variable). The
  variation in the values of the bounds of r is across pairs of
  variables. In other words, the bounds of r are specific for each pair.
\item
  rmin max can be -1, and rmax can be 1, but neither of these values
  implies the other (as seen in examples (b) and (c). The combination of
  rmin=-1 and rmax=1 is possible, but dependent on specific
  characteristics of the two marginal distributions involved (see
  examples (d) and (f). Both variables must not only have identical
  distributions, but these must also be symmetrical.
\item
  rmin and rmax are not necessarily symmetric around 0. Such symmetry
  may exist, but requires specific characteristics of the pair of
  marginal distributions, namely that both are symmetric (see examples
  (e), (g) and (h).\\
\item
  The extent to which the reference interval {[}-1, 1{]} is covered by
  the interval {[}rmin, rmax{]} varies across pairs of variables.
\item
  The values of rmin and rmax depend not only on the marginal
  distributions of the two variables, but also on the values assigned to
  their categories. Unless specified differently, we follow in this
  article the common practice of assigning consecutive integer values to
  the ordered categories. But any monotonic transformation of values is
  equally adequate to represent the ordinal character of variables.
  Examples (i) and (j) involve the same pair of variables, but in (i)
  their categories have been assigned values 1. 2, 3 and 4, while in (j)
  the assigned values were 1, 2.25, 2.75 and 4. The consequences of this
  change are visible in the values of rmin and rmax. This observation
  implies that value assignment to categories of orcats (and the
  justification thereof) remains an issue of importance when using
  Pearson correlations on orcats. It is a different problem, however,
  than that of constrainment of correlations on which we focus here.
\end{itemize}

Consequences of constrained correlations in applied research

\begin{itemize}
\tightlist
\item
  Descriptive interpretation of correlations (Cohen's criteria) is
  undermined
\item
  It is also undermined by every correlation being expressed on its own
  unique scale/unique calibration
\item
  Need a nice example of 2 correlations with \(r(a,b) \le r(b,c)\) and
  where \(r(a,b)\)max is \textless{} \(r(b.c)\): is it smaller because
  the linear relationship is weaker, or because he constrainment is
  stronger?
\end{itemize}

\subsection{Background and
Definitions}\label{background-and-definitions}

\begin{itemize}
\tightlist
\item
  Ordinal data and assumptions of interval-scale correlation\\
\item
  Existing alternatives: Spearman, polychoric, polyserial; their
  strengths and limitations\\
\item
  Coupling and rearrangement in discrete settings
\end{itemize}

\section{Theoretical Framework}\label{theoretical-framework}

The problem of extremal correlations under fixed marginals connects directly to classical results in probability inequalities. Fréchet’s work on admissible joint distributions (\autocite{frech51-sur}) and Hoeffding’s development of covariance bounds (\autocite{hoef40-massta}) laid the groundwork. Later refinements, such as Rüschendorf’s demonstration of sharpness of Fréchet bounds (\autocite{rues81-sharpn}), clarified when comonotonic and countermonotonic couplings achieve extremal values. The rearrangement inequality, and its discrete extensions, ensure that matching high values of \( X \) with high values of \( Y \) maximizes \(\mathbb{E}[XY]\), while opposite pairings minimize it. This perspective is elaborated in comparison and coupling frameworks (\autocite{muel02-compa}; \autocite{thor00-coupli}). These results provide the theoretical backbone for our derivations of \(r_{\min}\) and \(r_{\max}\).

\subsection{Optimal Coupling and Correlation
Bounds}\label{optimal-coupling-and-correlation-bounds}

We formalize the problem of identifying the joint distribution
(coupling) of two ordinal variables with fixed marginals that maximizes
or minimizes the Pearson correlation.

\phantomsection\label{prop-max-product}
\begin{proposition}

Let $X$ and $Y$ be discrete ordinal random variables, each with finite ordered support:

$$
X: x_1 < x_2 < \dots < x_{K_X},\quad Y: y_1 < y_2 < \dots < y_{K_Y}
$$

with marginal probability distributions $p_X(x_i) = P(X = x_i)$, $p_Y(y_j) = P(Y = y_j)$, respectively.

Then, the maximum and minimum values of the Pearson correlation coefficient $r_{XY}$, consistent with the fixed marginals, are obtained as follows:

- Maximum $r_{XY}$: achieved by pairing largest $X$-values with largest $Y$-values (comonotonic arrangement).
    
- Minimum $r_{XY}$: achieved by pairing largest $X$-values with smallest $Y$-values (anti-comonotonic arrangement).
\end{proposition}

\begin{proof}
The Pearson correlation coefficient is given by:

\[r_{XY} = \frac{E[XY] - E[X]E[Y]}{\sigma_X \sigma_Y}.\]

Given fixed marginals \(p_X(x_i)\), \(p_Y(y_j)\), the expectations
\(E[X], E[Y]\) and standard deviations \(\sigma_X, \sigma_Y\) are
constant. Thus, to maximize or minimize \(r_{XY}\), we only need to
maximize or minimize the expectation \(E[XY]\):

\[E[XY] = \sum_{i=1}^{K_X}\sum_{j=1}^{K_Y} x_i y_j P(X = x_i, Y = y_j).\]

This proof relies on a rearrangement inequality of
\autocite{hard52-inequalities} (Theorem 368):

For real sequences \(a_1 \leq a_2 \leq \dots \leq a_n\) and
\(b_1 \leq b_2 \leq \dots \leq b_n\), and for any permutation \(\pi\) of
indices \(\{1,\dots,n\}\) we have the following:

\[a_1 b_1 + a_2 b_2 + \dots + a_n b_n \geq a_1 b_{\pi(1)} + a_2 b_{\pi(2)} + \dots + a_n b_{\pi(n)}\]
and

\[a_1 b_n + a_2 b_{n-1} + \dots + a_n b_1 \leq a_1 b_{\pi(1)} + a_2 b_{\pi(2)} + \dots + a_n b_{\pi(n)}\]

Equality occurs only if the permutation \(\pi\) results in a
monotonically increasing sequence (for maximum) or monotonically
decreasing sequence (for minimum).

To explicitly connect this to our problem, enumerate observations
explicitly from sorted marginals:

\[
X_{\downarrow}: x_{[K_X]} \geq x_{[K_X-1]} \geq \dots \geq x_{[1]}, \quad  
        Y_{\downarrow}: y_{[K_Y]} \geq y_{[K_Y-1]} \geq \dots \geq y_{[1]}.
\] (Sorted in descending order for maximization). Next expand discrete
probability distributions into explicit observation-level sequences. If
\(K_X \neq K_Y\), extend the shorter sequence with zero-probability
categories so both sequences have equal length, which does not affect
marginal probabilities or correlation.

Hence, to maximize \(E[XY]\): pair the largest ranks of \(X\) with the
largest ranks of \(Y\). Explicitly, start from the top (highest values)
and move downward:

\begin{itemize}
\item
  Pair as many observations of the largest \(X\)-category as possible
  with the largest \(Y\)-category, then next-largest, etc., until all
  observations are paired.
\item
  By the Hardy--Littlewood--Pólya Rearrangement Inequality
  \autocite{hard52-inequalities}, this explicitly constructed pairing
  yields the maximum possible sum of products, hence maximizing
  \(E[XY]\).
\end{itemize}

Likewise, to minimize \(E[XY]\) pair the largest \(X\)-categories with
the smallest \(Y\)-categories, then second-largest \(X\)-categories with
second-smallest \(Y\)-categories, and so forth (anti-comonotonic
ordering). By the rearrangement inequality, this arrangement explicitly
yields the minimal sum of products, hence minimizing \(E[XY]\).
\end{proof}

In words, given fixed marginals, we have shown:

\begin{itemize}
\item
  \textbf{Maximum correlation} \(r_{\text{max}}\) is achieved by
  comonotonic (e.g.~sorting values of \(X\) and \(Y\) from largest to
  smallest independently and pairing) arrangement:

  \[P_{\text{max}}(X=x_{[i]}, Y=y_{[i]}) \quad\text{in either descending or ascending order.}\]
\item
  \textbf{Minimum correlation} \(r_{\text{min}}\) is achieved by
  anti-comonotonic arrangement:

  \[P_{\text{min}}(X=x_{[i]}, Y=y_{[n+1-i]}) \quad\text{pairing descending X with ascending Y, or viceversa.}\]
\end{itemize}

With this explicit construction we can simply calculate \(r_{max}\) and
\(r_{min}\) when the values of \(X\) and \(Y\) are ranks.

\phantomsection\label{corr-min}
\begin{corollary} {(Analytic Forms for $r_{\min}$ and $r_{\max}$)}


Let two ordinal variables $X$ and $Y$ have values as follows:

- $X$ with $K_X$ ordinal levels: $1,2,\dots,K_X$
    
- $Y$ with $K_Y$ ordinal levels: $1,2,\dots,K_Y$
    
with absolute frequencies:

- $n_{X}(i)$, number of observations at level $i$ of $X$.
    
- $n_{Y}(j)$, number of observations at level $j$ of $Y$.
    

Then the maximum Pearson correlation is:

$$r_{\text{max}} = \frac{E[XY]_{\text{max}} - \bar{X}\bar{Y}}{\sigma_X \sigma_Y}.$$
Where 

$$
E[XY]_{\max} =\frac{1}{N}\sum_{i=1}^{K_X}\sum_{j=1}^{K_Y}x_i y_j M_{i,j}
$$

where explicitly defined frequencies $M_{i,j}$ pair observations in a comonotonic manner, calculated explicitly via the greedy iterative algorithm in as Algorithm  \ref{alg:greedy} in the proof.


Then explicitly the minimal expectation is:

$$E[XY]_{\text{min}}=\frac{1}{N}\sum_{i=1}^{K_X}\sum_{j=1}^{K_Y} x_i y_j R_{i,j}$$
Where $R_{ij}$ is defined in the proof.
\end{corollary}

\begin{proof}
Algorithm for explicit calculation of \(M_{i,j}\) (greedy method):

Initialize available frequencies explicitly:
\(n'_Y(j)=n_Y(j)\quad\text{for each }j\)

\begin{itemize}
\item
  For each level \(x_i\), from smallest \(i=1\) to largest:

  \begin{itemize}
  \item
    For each level \(y_j\), from smallest \(j=1\) to largest:

    \begin{itemize}
    \tightlist
    \item
      Pair as many observations as possible:
    \end{itemize}
  \end{itemize}
\end{itemize}

\[M_{i,j}=\min\left(n_X(i)-\sum_{s=1}^{j-1}M_{i,s},\, n'_Y(j)\right)\]

Next, update explicitly: \[n'_Y(j)\leftarrow n'_Y(j)-M_{i,j}\]

Define the matrix \(M_{i,j}\), the frequency with which the level
\(x_i\) pairs with the level \(y_j\). Clearly, we must have:

\[\sum_j M_{i,j}=n_X(i),\quad\sum_i M_{i,j}=n_Y(j)\]

Similarly, the minimum Pearson correlation (\(r_{\text{min}}\)) achieved
by anti-comonotonic alignment (pairing largest ranks of \(X\) with
smallest ranks of \(Y\)) is explicitly:
\[r_{\text{min}} = \frac{E[XY]_{\text{min}} - \bar{X}\bar{Y}}{\sigma_X \sigma_Y}.\]

To achieve the maximum expected product \(E[XY]_{\text{max}}\), do the
following: First sort ranks from largest to smallest independently for
both \(X\) and \(Y\); then pair ranks directly from highest to lowest
available observations.

\textbf{Minimum Expectation}

The formula for \(E[XY]_{\text{min}}\) is derived in a similar manner.
Define explicitly a matrix \(R_{i,j}\) giving frequencies of pairings
\(X=x_i\) with \(Y=y_j\). For the minimal expectation, pair largest
available \(X\)-values explicitly with smallest available \(Y\)-values
first (greedy algorithm):

Initialize explicitly available frequencies:

\[n'_Y(j)=n_Y(j)\quad\text{for each }j\]

\begin{itemize}
\item
  For \(i\) from largest to smallest (\(i=K_X\) down to \(1\)):

  \begin{itemize}
  \tightlist
  \item
    For \(j\) from smallest to largest (\(j=1\) up to \(K_Y\)):

    \begin{itemize}
    \tightlist
    \item
      Pair explicitly as many observations as possible:
      \[R_{i,j}=\min\left(n_X(i)-\sum_{s=1}^{j-1}R_{i,s},\, n'_Y(j)\right)\]
    \end{itemize}
  \end{itemize}
\item
  Update explicitly available frequencies:
\end{itemize}

\[n'_Y(j)\leftarrow n'_Y(j)-R_{i,j}\]

Then explicitly the minimal expectation is:

\[E[XY]_{\text{min}}=\frac{1}{N}\sum_{i=1}^{K_X}\sum_{j=1}^{K_Y} x_i y_j R_{i,j}\]

This explicit equation applies Whitt's rearrangement theorem ,
\textcite{whit76-bivariate} , ensuring proper maximal pairing of ranks.
(\textcite{whit76-bivariate})
\end{proof}

\textbf{Discussion}

We want to construct a joint distribution table \(M\) such that:

\begin{itemize}
\item
  The \textbf{row sums} match the marginal counts for \(X\):
  \(\sum_j M_{i,j} = n_X(i)\)
\item
  The \textbf{column sums} match the marginal counts for \(Y\):
  \(\sum_i M_{i,j} = n_Y(j)\)
\end{itemize}

To do this, we go \textbf{greedily}, pairing the highest available \(X\)
values with the \textbf{smallest available} \(Y\) values, one cell at a
time.

To avoid \textbf{exceeding} the available marginal totals in \(Y\), we
must track how much of \(n_Y(j)\) is \textbf{still unassigned} at each
step --- and that's what \(n_Y'(j)\) represents.\\
After assigning \(M_{i,j}\) units to the cell \((i,j)\), we subtract
that amount from the remaining pool of level \(y_j\) values, like this:
\(n_Y'(j) \leftarrow n_Y'(j) - M_{i,j}\)This ensures we don't
\textbf{overfill} any column of the matrix.

To maximize the expectation \(E[XY]\), the rearrangement inequality
states we must pair values of \(X\) and \(Y\) from smallest-to-smallest,
second-smallest-to-second-smallest, and so forth (comonotonic
alignment).

Thus, explicitly:

\begin{itemize}
\tightlist
\item
  Sort \(X\): lowest-to-highest:
\end{itemize}

\[X_{\text{sorted}} = (\underbrace{x_1,\dots,x_1}_{n_X(1)}, \underbrace{x_2,\dots,x_2}_{n_X(2)}, \dots,\underbrace{x_{K_X},\dots,x_{K_X}}_{n_X(K_X)})\]

Sorted \(Y\)-values (ascending order for max, descending for min):

\[Y_{\text{sorted(max)}} = (\underbrace{y_1,\dots,y_1}_{n_Y(1)}, \underbrace{y_2,\dots,y_2}_{n_Y(2)}, \dots,\underbrace{y_{K_Y},\dots,y_{K_Y}}_{n_Y(K_Y)})\]

\[Y_{\text{sorted(min)}} = (\underbrace{y_{K_Y},\dots,y_{K_Y}}_{n_Y(K_Y)}, \underbrace{y_{K_Y-1},\dots,y_{K_Y-1}}_{n_Y(K_Y-1)}, \dots,\underbrace{y_1,\dots,y_1}_{n_Y(1)})\]

The maximum and minimum expectations are simply:

\textbf{Maximum expectation} (pair element-by-element):

\[E[XY]_{\text{max}} = \frac{1}{N}\sum_{k=1}^{N} X_{\text{sorted}}(k)\cdot Y_{\text{sorted(max)}}(k)\]

\begin{itemize}
\tightlist
\item
  \textbf{Minimum expectation}:
\end{itemize}

\[E[XY]_{\text{min}} = \frac{1}{N}\sum_{k=1}^{N} X_{\text{sorted}}(k)\cdot Y_{\text{sorted(min)}}(k)\]

These classical bounds define feasible joint distributions given
marginals and constrain the attainable values of \(r\). There is an
interesting connection to Fréchet--Hoeffding and Boole--Fréchet
Inequalities, which we note but leave for future research.\footnote{Fréchet--Hoeffding
  and Boole--Fréchet Inequalities: The upper bound is always attained by
  \(r_{max}\). However, the lower bound is \emph{not} guaerenteed to
  attain. The fundamental asymmetry arises because discrete ordinal
  distributions cannot smoothly bridge these gaps. Thus, the upper bound
  alignment is always feasible, while the lower bound alignment
  frequently encounters structural obstacles.

  In the discrete ordinal setting, \(r_{\max}\) is \textbf{always}
  attained by the FH upper (comonotonic) construction.\\
  \emph{(BFbounds.md; Proofs-OrdinalVariables.md; Ordinal Variables and
  Correlation.md)}

  The FH lower (countermonotonic) construction may be
  \textbf{unattainable} in discrete settings; \(r_{\min}\) is then the
  best feasible opposite-order arrangement.\\
  \emph{(BFbounds.md; Proofs-OrdinalVariables.md; Ordinal Variables and
  Correlation.md)}

  \textbf{Upper bound} (positive alignment):

  \begin{itemize}
  \item
    Sorting \(X\) and \(Y\) separately in descending order, then pairing
    largest-to-largest, smallest-to-smallest, \textbf{always works}.
  \item
    You never run out of ``matching'' observations because the
    descending order for both variables naturally aligns category sizes.
    Even if one variable has large gaps, you can still sequentially pair
    without difficulty.
  \end{itemize}

  \textbf{Lower bound} (negative alignment):

  \begin{itemize}
  \item
    You sort \(X\) descending but \(Y\) ascending. You're pairing
    extremes together: highest \(X\) with lowest \(Y\), lowest \(X\)
    with highest \(Y\).
  \item
    Now you face potential ``imbalance'':

    \begin{itemize}
    \item
      If one category is large for \(X\) but the corresponding
      opposite-ranked category for \(Y\) is small, you quickly exhaust
      available ``opposite-end'' observations.
    \item
      After exhausting the best opposites, you are forced to pair less
      ``extreme'' pairs, creating a suboptimal alignment.
    \end{itemize}
  \item
    Hence, perfect negative alignment (lower bound) becomes impossible
    to achieve fully, leaving you short of the theoretical
    Fréchet--Hoeffding lower bound.
  \end{itemize}}

\subsubsection{Illustration}\label{illustration}

To illustrate these idea concretely, Table 1 presents an illustrative
example of comonotonic and anti-comonotonic arrangements and of the
resulting lower- and upper-bounds of r, for two 4-category variables and
12 cases. One can easily expand this example to pairs of variables with
different numbers of categories and different numbers of cases.

% Table 1: Determining minimum and maximum attainable values
\begin{table}[H]
\centering
\caption{Determining minimum and maximum attainable values of correlation between two ordered categorical variables X and Y given their marginal distributions.}
\label{tab:correlation_bounds}

\vspace{0.5em}

% Part (a): Observed frequency distributions
\textbf{(a) Observed frequency distributions}

\vspace{0.3em}

\begin{tabular}{c|cccc|c}
\hline
 & \textbf{Category 1} & \textbf{Category 2} & \textbf{Category 3} & \textbf{Category 4} & \textbf{Total} \\
\hline
\textbf{X} & 6 & 3 & 2 & 1 & 12 \\
\textbf{Y} & 3 & 5 & 2 & 2 & 12 \\
\hline
\end{tabular}

\vspace{0.5em}

% Part (b): Comonotonic arrangement
\textbf{(b) Comonotonic arrangement and pairing:}

\vspace{0.3em}

\begin{tabular}{ll}
\textbf{X:} & 1 1 1 1 1 1 2 2 2 3 3 4 \\
\textbf{Y:} & 1 1 1 2 2 2 2 2 3 3 4 4 \quad $r_{\max} = 0.87831$
\end{tabular}

\vspace{0.5em}

% Part (c): Anti-comonotonic arrangement
\textbf{(c) Anti-comonotonic arrangement and pairing:}

\vspace{0.3em}

\begin{tabular}{ll}
\textbf{X:} & 1 1 1 1 1 1 2 2 2 3 3 4 \\
\textbf{Y:} & 4 4 3 3 2 2 2 2 2 1 1 1 \quad $r_{\min} = -0.79466$
\end{tabular}

\end{table}

% Table 2: Illustrative pairs of marginal distributions
\begin{table}[H]
\centering
\caption{Illustrative pairs of marginal distributions and their associated minimum and maximum attainable values of $r$ ($r_{\min}$, $r_{\max}$) given these distributions. The distributions are given in counts (n=12). The categories are assigned consecutive integer values starting with 1. C1 and C2 are measures of constraint, as discussed in the text.}
\label{tab:marginal_distributions}

\vspace{0.5em}

\begin{tabular}{c|cccc|c|cc}
\hline
 & \multicolumn{4}{c|}{\textbf{Category}} & \textbf{Total} & $r_{\min}$ & $r_{\max}$ \\
 & \textbf{1} & \textbf{2} & \textbf{3} & \textbf{4} & & & \\
\hline
\multicolumn{8}{l}{\textbf{(a) Two arbitrarily distributed ordinal categorical variables (orcats)}} \\
\textbf{X} & 6 & 3 & 2 & 1 & 12 & -0.79466 & 0.87831 \\
\textbf{Y} & 3 & 5 & 2 & 2 & 12 & & \\
\hline
\multicolumn{8}{l}{\textbf{(b) One arbitrarily distributed variable, and one with same distribution inversed}} \\
\textbf{X} & 6 & 3 & 2 & 1 & 12 & -1 & 0.71429 \\
\textbf{Y} & 1 & 2 & 3 & 6 & 12 & & \\
\hline
\multicolumn{8}{l}{\textbf{(c) Two variables with arbitrary identical distributions}} \\
\textbf{X} & 6 & 3 & 2 & 1 & 12 & -0.71429 & 1 \\
\textbf{Y} & 6 & 3 & 2 & 1 & 12 & & \\
\hline
\multicolumn{8}{l}{\textbf{(d) Two variables with identical symmetric and unimodal distributions}} \\
\textbf{X} & 2 & 4 & 4 & 2 & 12 & -1 & 1 \\
\textbf{Y} & 2 & 4 & 4 & 2 & 12 & & \\
\hline
\multicolumn{8}{l}{\textbf{(e) Two variables with different distributions, both symmetric and unimodal}} \\
\textbf{X} & 2 & 4 & 4 & 2 & 12 & -0.91169 & 0.91169 \\
\textbf{Y} & 1 & 5 & 5 & 1 & 12 & & \\
\hline
\multicolumn{8}{l}{\textbf{(f) Two variables with identical symmetric bimodal distributions}} \\
\textbf{X} & 4 & 2 & 2 & 4 & 12 & -1 & 1 \\
\textbf{Y} & 4 & 2 & 2 & 4 & 12 & & \\
\hline
\multicolumn{8}{l}{\textbf{(g) Two variables with different distributions, both symmetric and bimodal}} \\
\textbf{X} & 4 & 2 & 2 & 4 & 12 & -0.95673 & 0.95673 \\
\textbf{Y} & 5 & 1 & 1 & 5 & 12 & & \\
\hline
\multicolumn{8}{l}{\textbf{(h) One symmetric unimodal variable, the other symmetric bimodal}} \\
\textbf{X} & 2 & 4 & 4 & 2 & 12 & -0.89923 & 0.89923 \\
\textbf{Y} & 4 & 2 & 2 & 4 & 12 & & \\
\hline
\multicolumn{8}{l}{\textbf{(i) Two variables, one very left skewed, the other very right skewed}} \\
\textbf{X} & 1 & 1 & 1 & 9 & 12 & -0.95818 & 0.31939 \\
\textbf{Y} & 8 & 2 & 1 & 1 & 12 & & \\
\hline
\multicolumn{8}{l}{\textbf{(j) Same as situation (i), above, but categories assigned values 1; 2.25; 2.75; 4}} \\
\textbf{X} & 1 & 1 & 1 & 9 & 12 & -0.93321 & 0.33829 \\
\textbf{Y} & 8 & 2 & 1 & 1 & 12 & & \\
\hline
\end{tabular}

\end{table}

\subsection{Examples}\label{examples}

Below, ``scores \(0{:}(K-1)\)'' means category values
\([0,1,\dots,K-1]\). ``Affine recode'' means
\(a+b\times\text{(old score)}\) with \(b>0\). ``Non-affine monotone''
means strictly increasing but not affine (e.g., convex spacing).

\subsubsection{Notation (permutation extremes, ``Lane
A'')}\label{notation-permutation-extremes-lane-a}

\begin{itemize}
\item
  Build expanded vectors \(x\) and \(y\) from \textbf{counts} and
  \textbf{scores}.
\item
  \(r_{\max}=\mathrm{cor}(\text{sort}(x,\downarrow),\text{sort}(y,\downarrow))\).
\item
  \(r_{\min}=\mathrm{cor}(\text{sort}(x,\downarrow),\text{sort}(y,\uparrow))\).
\end{itemize}

\begin{center}\rule{0.5\linewidth}{0.5pt}\end{center}

\subsubsection{\texorpdfstring{A) \textbf{Same number of categories, but
\(|r_{\min}|\ne r_{\max}\)} (concrete
numbers)}{A) Same number of categories, but \textbar r\_\{\textbackslash min\}\textbar\textbackslash ne r\_\{\textbackslash max\} (concrete numbers)}}\label{a-same-number-of-categories-but-r_minne-r_max-concrete-numbers}

\textbf{Counts (both \[K=5\]):}

\begin{itemize}
\item
  \[X:\ [100,\,200,\,300,\,200,\,200]\]
\item
  \(Y:\ [100,\,100,\,100,\,300,\,400]\) (skewed; not palindromic)
\end{itemize}

\textbf{Scores:} \(0{:}4\) for both.

\begin{itemize}
\item
  \[r_{\max}\approx 0.929398758\]
\item
  \[r_{\min}\approx -0.881118303\]
\item
  \[|r_{\min}| \approx 0.8811 \ne r_{\max}\]
\end{itemize}

\textbf{Affine recode (10,20,30,40,50) for both:}

\begin{itemize}
\tightlist
\item
  \(r_{\max}, r_{\min}\) \textbf{unchanged} (up to rounding).
\end{itemize}

\textbf{Non-affine monotone recode for \(Y\):} \([0,1,1.5,3,10]\).

\begin{itemize}
\item
  \[r_{\max}\approx 0.930261031\]
\item
  \[r_{\min}\approx -0.863250703\]
\item
  Values change (Pearson is \textbf{not} invariant to non-affine
  spacing); magnitudes still unequal.
\end{itemize}

🧠 \textbf{Why unequal magnitudes?} No palindromic symmetry in \(Y\) ⇒
the ``palindrome-sum'' pattern in \(y'-\)space isn't zero, so
\(S_{\max}+S_{\min}\ne 0\).

\begin{center}\rule{0.5\linewidth}{0.5pt}\end{center}

\subsubsection{\texorpdfstring{B) \textbf{Different numbers of
categories, yet \(|r_{\min}|=r_{\max}\)} (concrete
numbers)}{B) Different numbers of categories, yet \textbar r\_\{\textbackslash min\}\textbar=r\_\{\textbackslash max\} (concrete numbers)}}\label{b-different-numbers-of-categories-yet-r_minr_max-concrete-numbers}

\textbf{Counts:}

\begin{itemize}
\item
  \(X\) (4 categories): \([100,\,200,\,300,\,400]\)
\item
  \(Y\) (5 categories): \([100,\,200,\,400,\,200,\,100]\)
  (\textbf{palindromic})
\end{itemize}

\textbf{Scores:} \(X:\ 0{:}3\), \(Y:\ 0{:}4\).

\begin{itemize}
\item
  \[r_{\max}\approx 0.912870929\]
\item
  \[r_{\min}\approx -0.912870929\]
\item
  \(|r_{\min}| = r_{\max}\) (to machine precision)
\end{itemize}

\textbf{Affine recode (e.g., 10,20,\ldots):}

\begin{itemize}
\tightlist
\item
  Magnitudes remain \textbf{identical}; affine transforms don't change
  Pearson correlation.
\end{itemize}

\textbf{Non-affine monotone recode for \(Y\):} \([0,1,1.5,3,10]\).

\begin{itemize}
\item
  \[r_{\max}\approx 0.580072921\]
\item
  \[r_{\min}\approx -0.842041337\]
\item
  \textbf{Magnitudes no longer match} because antisymmetry of the
  \textbf{scores} about the center is lost, even though the
  \textbf{counts} are palindromic.
\end{itemize}

🧠 \textbf{Why equal magnitudes in the affine case?} Palindromic
\textbf{mass} symmetry in \(Y\) + \textbf{antisymmetric scoring} (any
affine map of \(0{:}(K_Y-1)\)) ⇒ the palindrome-sum vector
\(y'_i+y'_{n+1-i}\equiv 0\). Hence \(S_{\max}+S_{\min}=0\) ⇒
\(|r_{\min}|=r_{\max}\) for \textbf{any} \(X\), irrespective of
\(K_X\neq K_Y\).

\begin{center}\rule{0.5\linewidth}{0.5pt}\end{center}

\subsubsection{\texorpdfstring{C) \textbf{Spearman invariance (same data
as in
B)}}{C) Spearman invariance (same data as in B)}}\label{c-spearman-invariance-same-data-as-in-b}

Compute ranks (with average ties) \textbf{once} for \(x\) and \(y\),
then form extremes by sorting those rank vectors.

\begin{itemize}
\item
  With \(Y\) scored \(0{:}4\):\\
  \(r_{\max}^{\text{Spearman}}\approx 0.920837215,\ \ r_{\min}^{\text{Spearman}}\approx -0.920837215\).
\item
  With \(Y\) scored non-affinely \([0,1,1.5,3,10]\):\\
  \textbf{identical} Spearman extremes: \(0.920837215\) and
  \(-0.920837215\).
\end{itemize}

✅ \textbf{Spearman is invariant} to any strictly monotone recode
because ranks depend only on order, not spacing.

\begin{center}\rule{0.5\linewidth}{0.5pt}\end{center}

\subsection{what changes and what doesn't (Pearson vs
Spearman)}\label{what-changes-and-what-doesnt-pearson-vs-spearman}

\begin{itemize}
\item
  \textbf{Pearson:}

  \begin{itemize}
  \item
    \textbf{Affine} recodes \(a+b\cdot\text{score}\) (with \(b>0\))
    leave \(r_{\max}, r_{\min}\), and any \(|r_{\min}|=r_{\max}\)
    conclusions \textbf{unchanged}.
  \item
    \textbf{Non-affine monotone} recodes can change
    \(r_{\max}, r_{\min}\) and can break \(|r_{\min}|=r_{\max}\) even
    with palindromic counts.
  \end{itemize}
\item
  \textbf{Spearman:}

  \begin{itemize}
  \item
    Any strictly monotone recode leaves extremes \textbf{unchanged}.
  \item
    Palindromic counts in \(Y\) (with the usual symmetric rank scheme)
    give \(|r_{\min}|=r_{\max}\).
  \end{itemize}
\end{itemize}

\subsection{Population Bounds}\label{population-bounds}

In this section we further explore the theoretical properties on
\(r_{min}\) and \(r_{max}\).

\subsubsection{Lemma (signs of extreme correlations under monotone
scoring)}\label{lemma-signs-of-extreme-correlations-under-monotone-scoring}

\phantomsection\label{prop-max-product}
\begin{lemma}

Let $X$ and $Y$ be ordinal variables. Assign numeric scores to categories via **non-decreasing** maps $s_X, s_Y$ that respect the category order (scores may be affine or nonlinear; ties allowed). Define two extreme couplings:

- **Comonotone (upper)**: pair larger scored values of $X$ with larger scored values of $Y$.
- **Antitone (lower)**: pair larger scored values of $X$ with smaller scored values of $Y$.
    

Let $r_{\max}$ and $r_{\min}$ be the Pearson correlations computed from these couplings, either (i) in the **finite-sample/permutation** sense (sort scored observations in the same vs. opposite order), or (ii) in the **Fréchet–Hoeffding (FH)/population** sense (pair quantiles $Q_X(u)$ with $Q_Y(u)$ vs. $Q_Y(1-u)$).

Then:

1. $r_{\max}\ge 0$ and $r_{\min}\le 0$.
    
2. Moreover, $r_{\max}>0$ and $r_{\min}<0$ **iff** both variables have **positive variance** under the chosen scoring (i.e., neither scored variable is almost surely constant).
    

_(If you reverse the scoring of one variable (decreasing map), swap “upper” and “lower”. For Spearman’s $$\rho$$, take the scores to be midranks—monotone by construction—and the same conclusions hold.)_
\end{lemma}

\begin{proof}
Proof

We give parallel proofs in the finite-sample and population settings.

\begin{enumerate}
\def\labelenumi{\Alph{enumi})}
\tightlist
\item
  Finite-sample (permutation) setting
\end{enumerate}

Let \(x_{(1)}\ge\cdots\ge x_{(n)}\) and \(y_{(1)}\ge\cdots\ge y_{(n)}\)
be the \textbf{scored} observations of \(X\) and \(Y\) sorted in
nonincreasing order (monotone scoring ensures sorting by category equals
sorting by score). Center them:

\[x'_i := x_{(i)}-\bar x,\qquad y'_i := y_{(i)}-\bar y,\quad  
\bar x := \tfrac1n\sum_i x_{(i)},\;\bar y := \tfrac1n\sum_i y_{(i)}.\]

Define the two extreme centered sums

\[S_{\max} := \sum_{i=1}^n x'_i\,y'_i,\qquad  
S_{\min} := \sum_{i=1}^n x'_i\,y'_{\,n+1-i}.\]

\begin{itemize}
\item
  \textbf{Signs.}\\
  Chebyshev's sum inequality (rearrangement for similarly/ oppositely
  ordered sequences) gives

  \[\frac1n\sum_i x'_i y'_i \;\ge\; \Big(\tfrac1n\sum_i x'_i\Big)\Big(\tfrac1n\sum_i y'_i\Big)=0,  
    \quad\Rightarrow\; S_{\max}\ge 0,\]

  and, because \(y'_{\,n+1-i}\) is \textbf{nondecreasing} in \(i\),

  \[\frac1n\sum_i x'_i y'_{\,n+1-i} \;\le\; \Big(\tfrac1n\sum_i x'_i\Big)\Big(\tfrac1n\sum_i y'_{\,n+1-i}\Big)=0,  
    \quad\Rightarrow\; S_{\min}\le 0.\]

  Dividing by \(\sigma_X\sigma_Y>0\) (the standard deviations of the
  scored variables) yields \(r_{\max}\ge 0\) and \(r_{\min}\le 0\).
\item
  \textbf{Strictness \((\Leftrightarrow\) positive variances).}\\
  Suppose \(\sigma_X>0\) and \(\sigma_Y>0\) (i.e., neither \(\{x'_i\}\)
  nor \(\{y'_i\}\) is identically zero). The rearrangement inequality
  says \(S_{\max}\) is the \textbf{maximum} and \(S_{\min}\) the
  \textbf{minimum} of \(\sum_i x'_i y'_{\pi(i)}\) over all permutations
  \(\pi\); moreover \(S_{\min}<S_{\max}\). The \textbf{average} over all
  permutations equals

  \[\frac1{n!}\sum_{\pi}\sum_i x'_i y'_{\pi(i)}  
    = \sum_i x'_i \Big(\frac1{n!}\sum_{\pi} y'_{\pi(i)}\Big)  
    = \sum_i x'_i \cdot \frac1n\sum_j y'_j  
    = 0.\]

  Since the minimum is \(<\) the maximum and the average is \(0\), it
  follows that \(S_{\min}<0<S_{\max}\). Hence \(r_{\min}<0<r_{\max}\).

  Conversely, if \(\sigma_X=0\) or \(\sigma_Y=0\), then one centered
  sequence is identically zero, so \(S_{\max}=S_{\min}=0\) and the
  correlations are undefined or (in a limiting sense) attain 0. This
  matches the ``only if'' direction.
\end{itemize}

\begin{enumerate}
\def\labelenumi{\Alph{enumi})}
\setcounter{enumi}{1}
\tightlist
\item
  Population (FH/quantile) setting
\end{enumerate}

Let \(Q_X, Q_Y : [0,1]\to\mathbb{R}\) be the quantile functions of the
\textbf{scored} variables. Set

\[f(u):=Q_X(u)-\mu_X,\qquad g(u):=Q_Y(u)-\mu_Y,\]

so \(\int_0^1 f=\int_0^1 g=0\). Monotone scoring preserves category
order, hence \(Q_X, Q_Y\) are nondecreasing; therefore \(f, g\) are
nondecreasing functions.

\begin{itemize}
\item
  \textbf{Upper/comonotone covariance}:

  \[\mathrm{Cov}_{\max}=\int_0^1 f(u)\,g(u)\,du \;\ge\; \Big(\int_0^1 f\Big)\Big(\int_0^1 g\Big)=0\]

  by Chebyshev's \textbf{integral} inequality for similarly ordered
  functions.
\item
  \textbf{Lower/antitone covariance}:

  \[\mathrm{Cov}_{\min}=\int_0^1 f(u)\,g(1-u)\,du \;\le\; \Big(\int_0^1 f\Big)\Big(\int_0^1 g(1-u)\Big)=0,\]

  since \(f\) is nondecreasing and \(g(1-u)\) is \textbf{nonincreasing}.
\end{itemize}

Dividing by \(\sigma_X\sigma_Y>0\) yields \(r_{\max}\ge 0\) and
\(r_{\min}\le 0\).

\begin{itemize}
\tightlist
\item
  \textbf{Strictness.}\\
  If both variables have positive variance, then \(f\) and \(g\) are not
  a.e. constant. For nontrivial nondecreasing \(f,g\), Chebyshev's
  integral inequalities above are \textbf{strict} unless one function is
  a.e. constant. Hence \(\mathrm{Cov}_{\max}>0\) and
  \(\mathrm{Cov}_{\min}<0\), giving \(r_{\max}>0\) and \(r_{\min}<0\).
  Conversely, if either variance is \(0\), both covariances are \(0\).
\end{itemize}

This completes the proof. \(\square\)
\end{proof}

\subsection{Asymetry}\label{asymetry}

\textbf{Symmetry and the Role of Ties}

\begin{itemize}
\tightlist
\item
  Conditions under which \(r_{\min} = -r_{\max}\)
\end{itemize}

PROP: If one marginal is \textbf{palindromically symmetric} about its
center rank, then \(r_{\min} = -r_{\max}\).\\
\emph{(Cees -- r\_min MAGNITUDE Asymmetry in Pearson Correlation.md)}

Necessary/sufficient symmetry condition (precisely stated)

\begin{itemize}
\tightlist
\item
  \textbf{Pair-specific equality} \(|r_{\min}|=r_{\max}\) \(\iff\)
  \(\sum_i x'_i\,(y'_i+y'_{n+1-i})=0\).\\
  (Can happen by coincidence even if \(Y\) isn't symmetric.)
\item
  \textbf{Universal in \(X\)} equality (for \textbf{all} \(X\)):\\
  \textbf{Necessary and sufficient} that \textbf{one} marginal (say
  \(Y\)) has

  \begin{enumerate}
  \def\labelenumi{(\roman{enumi})}
  \tightlist
  \item
    \textbf{palindromic mass symmetry} \(p_Y(k)=p_Y(K_Y-1-k)\) and\\
  \item
    \textbf{antisymmetric scoring} \(s_Y(k)+s_Y(K_Y-1-k)=\text{const}\)
    (true for any affine map of \(0{:}(K_Y-1)\)).\\
    This condition is \textbf{independent of \(K_X\) vs \(K_Y\)}.
  \end{enumerate}
\end{itemize}

\subsubsection{1) setup (permutation extremes for a single
sample)}\label{setup-permutation-extremes-for-a-single-sample}

Let \(x_{(1)}\ge\cdots\ge x_{(n)}\) and \(y_{(1)}\ge\cdots\ge y_{(n)}\)
be the sorted sample values (or expanded from counts). Define centered
sequences

\[x'_i:=x_{(i)}-\bar x,\qquad y'_i:=y_{(i)}-\bar y,\]

and the two permutation extremes

\[S_{\max}=\sum_{i=1}^n x'_i\,y'_i,\qquad  
S_{\min}=\sum_{i=1}^n x'_i\,y'_{n+1-i},\]

with \(r_{\max}=S_{\max}/(\sigma_X\sigma_Y)\),
\(r_{\min}=S_{\min}/(\sigma_X\sigma_Y)\).

We always have \(S_{\max}\ge 0\) and \(S_{\min}\le 0\). Magnitude
symmetry \(|r_{\min}|=r_{\max}\) is equivalent to

\[S_{\max}+S_{\min}=0.\]

\subsubsection{\texorpdfstring{2) \textbf{pair-specific necessary \&
sufficient
condition}}{2) pair-specific necessary \& sufficient condition}}\label{pair-specific-necessary-sufficient-condition}

Let

\[z_i := y'_i + y'_{n+1-i}.\]

Then

\[S_{\max}+S_{\min}=\sum_{i=1}^n x'_i\,z_i.\]

Hence, for this \emph{particular} pair \((x,y)\),

\[\boxed{\ |r_{\min}|=r_{\max}\ \Longleftrightarrow\ \sum_i x'_i\,z_i=0\ } \tag{★}\]

(i.e., the centered \(X\) sequence is orthogonal to the
``palindrome-sum'' vector \(z\)).\\
This can hold even if \(Y\) is not symmetric---by coincidence of that
inner product being zero.

\subsubsection{\texorpdfstring{3) \textbf{universal condition
(independent of
\(X\))}}{3) universal condition (independent of X)}}\label{universal-condition-independent-of-x}

If you want \(|r_{\min}|=r_{\max}\) to hold for \textbf{every} possible
\(X\) (every centered, sorted \(x'\)), then (★) must hold for all
\(x'\). The only way is

\[z_i \equiv 0\quad\text{for all }i,\]

i.e.

\[y'_i + y'_{n+1-i}\equiv 0\ \ \ \forall i.\]

Translating back to categories and scores, a \textbf{necessary and
sufficient universal condition} is:

\begin{itemize}
\item
  \textbf{Palindromic mass symmetry for \(Y\):} \(p_Y(k)=p_Y(K_Y-1-k)\)
  for all \(k\), and
\item
  \textbf{Antisymmetric scoring for \(Y\):} the score map satisfies
  \(s_Y(k)+s_Y(K_Y-1-k)\) is constant (true for affine scores like
  \(0,1,\dots,K_Y-1\), and for Spearman midranks under palindromic
  counts).
\end{itemize}

Under these, \(z_i\equiv 0\) and thus \(|r_{\min}|=r_{\max}\) for any
\(X\).

\begin{quote}
If either mass symmetry or score antisymmetry fails, you can generally
find an \(X\) for which \(|r_{\min}|\ne r_{\max}\) (often
\(|r_{\min}|>r_{\max}\) when \(Y\) is skewed).
\end{quote}

\begin{itemize}
\tightlist
\item
  How ties in marginal distributions affect the attainability of the
  lower bound
\end{itemize}

\begin{enumerate}
\def\labelenumi{\arabic{enumi}.}
\tightlist
\item
  \textbf{\(r_{\max}\) = FH upper (comonotonic), always attainable} in
  discrete ordinal settings.
\item
  \textbf{\(r_{\min}\)} may \textbf{not} equal FH lower in discrete
  settings; use best feasible opposite-order arrangement.
\item
  No universal bound on relative magnitudes: \(|r_{\min}|\) may be less
  than, equal to, or greater than \(r_{\max}\).
\item
  Palindromic symmetry of one marginal guarantees
  \(|r_{\min}| = r_{\max}|\).
\end{enumerate}

VH 3. \textbf{Symmetry Analysis:}

\begin{itemize}
\tightlist
\item
  (3a) Visualize and explore how changes from symmetric to asymmetric
  marginal distributions influence the joint distributions and,
  consequently, the correlation. Consider the potential role of entropy
  in measuring asymmetry.
\end{itemize}

\begin{enumerate}
\def\labelenumi{\arabic{enumi}.}
\setcounter{enumi}{3}
\tightlist
\item
  \textbf{Investigate when the minimum and maximum correlations are
  asymmetric (\(r_{\text{min}} \neq -r_{\text{max}}\))}:
\end{enumerate}

\begin{itemize}
\tightlist
\item
  Show why and under what conditions this asymmetry occurs.
\item
  Visually investigate why perfect negative and positive correlations
  may not be symmetric in magnitude given ordinal marginals.
\end{itemize}

\textbf{Special Cases and Boundary Conditions}

\begin{itemize}
\tightlist
\item
  Equal vs.~unequal numbers of categories\\
\item
  Symmetric vs.~asymmetric marginal distributions
\end{itemize}

Often \(|r_{\min}| < r_{\max}\) for discrete asymmetrical marginals.\\
\emph{(BFbounds.md; Ordinal Variables and Correlation.md)} - \textbf{New
evidence (BES data)}: Possible to have \(|r_{\min}| > r_{\max}|\).\\
Example: \(X\) with 4 categories, \(Y\) with 5 categories, skewed
marginals;\\
\(r_{\max} = 0.85027\), \(r_{\min} = -0.89916\).\\
\emph{(Cees -- r\_min MAGNITUDE Asymmetry in Pearson Correlation.md)}

\textbf{Reason:} Pearson correlation uses centered values; skewness can
make negative deviations in antitone pairing larger in magnitude than
positive deviations in comonotone pairing.

VH

\subsection{Estiimation and Joint
Uncertainty}\label{estiimation-and-joint-uncertainty}

VH

\section{Measurement Issues and Metrics of
Constraint}\label{measurement-issues-and-metrics-of-constraint}

IN THIS SECTION: ``How to define metrics (on \emph{pairs} of discrete
probability distributions) to \emph{explain}:

VH

\begin{enumerate}
\def\labelenumi{\arabic{enumi}.}
\tightlist
\item
  \[\frac{ r_{max}-r_{min}}{2}\]
\item
  \[1-r_{max}\]
\item
  \[|-1-r_{min}|\]
\end{enumerate}

Or analogs for \(\hat{r}^2\):

\begin{enumerate}
\def\labelenumi{\arabic{enumi}.}
\tightlist
\item
  \[\frac{ r_{max}^2-r_{min}^2}{2}\] 1a.
  \[\frac{ (r_{max}-r_{min})^2}{2}\]
\item
  \[1-r_{max}^2\]
\item
  \[1-r_{min}^2\]
\end{enumerate}

As demonstrated above, the extent to which marginal distributions of
orcats constrain the range of attainable value of r varies considerably.
In the few illustrative examples presented in Table 2 rmin varied
between -1 and -0.71, while rmax varied between 0.32 and 1. Moreover, in
Table 2 rmin and rmax do not vary in lockstep but at least to some
extent independent of each other.{[}fn: Table 2 is a poor basis to
assess the functional relationship between rmin and rmax as it contains
a small number of handpicked situations. This relationship will be
investigated in detail later, in section ??.{]} This raises the question
how strong the overall constraining effect is that marginal
distributions exert on the attainable values of r. It also raises the
question of what factors affect the strength of this overall
constraining effect. In this section we describe two ways in which the
strength of constraints can be defined. Later (in section ??) we turn to
the questions of the conditions that drive the strength of this effect.
Perhaps the most obvious way to express the strength of the marginal
constraints on the attainable values of r is by comparing the width of
the interval between rmin and rmax with that of the ubiquitously used
reference interval {[}-1, 1{]}. As the width of the latter interval is
2, it is advantageous to express the strength of the overall
constraining effect --to be denoted as \(C_1\) as

\[
C_1 = \frac{|r_{min}| +|r_{max}|}{2}
\] \(C_1\) reflects the proportion of the range of the attainable values
of r compared to the theoretically maximum width of this range. Applying
this measure to the illustrative cases listed in Table 2 shows that C1
varies. Its highest value is 1, reflecting the absence of any
constraining effect (see examples (d) and (f) in Table 2). The lowest
value of C1 in Table 2 is for example (i), where it is 0.639. In other
words, the marginal distributions of the two variables in that example
restrict the range of its attainable values to 63.9\% of the traditional
{[}-1, 1{]} reference interval. Some might object to the \(C_1\) measure
on the grounds that it assumes that r (and thus also rmin and rmax) is
expressed at interval level and that one could thus state validly that a
correlation of 0.4 reflects an association twice as strong as one of
0.2. That would be so if numerical differences between correlations
correspond to similar differences in observable aspects of data
{[}insert references to correspondence NRS to ERS{]}. That not being the
case, one could instead compare r2s --percentages of shared variance--
rather than rs. {[}fn: This perspective is not uncontested, however, and
its usefulness depends on the purpose of inspecting correlations: as
reflections of the association between variables, or as reflections of
associations between indicators of a latent concept (for an elaboration
of this difference, cf.~Ozer 1985).{]} This would lead to a somewhat
different measure of the strength of the constraining effect of marginal
distributions on the attainable values of r, as follows:

\[
C_2 = \frac{r_{min}^2 +r_{max}^2}{2}
\]

For the examples in Table 2, the values C2 are, respectively, 0.701;
0.755; 0.715; 1; 0.831; 1; 0.809; 0.915 and 0.510. Obviously, C1 and C2
are correlated, but how strongly cannot be assessed very well from the
few handpicked examples in Table 2, but it will be considered below,
using more encompassing sets of data.

\subsection{Possible ``Metrics of
Constraint''}\label{possible-metrics-of-constraint}

\subsubsection{Overlap-Based Measures of
Asymmetry}\label{overlap-based-measures-of-asymmetry}

In addition to Total Variation (TV) distance, other metrics can quantify
the structural dissimilarity between a marginal distribution and its
reverse. One natural approach is to assess the degree of
\textbf{overlap} between the original distribution \(p\) and its reverse
\(p^{\mathrm{rev}}\), defined by reversing the order of the categories.

We consider three complementary measures of overlap and their
complements:

\paragraph{(i) Total Variation (TV)
Distance}\label{i-total-variation-tv-distance}

The TV distance, already introduced in Section 5.5, quantifies the
maximum absolute discrepancy between the two distributions:

\[
\text{TV}(p, p^{\mathrm{rev}}) = \frac{1}{2} \sum_{i=1}^{K} \left| p_i - p_{K+1 - i} \right|.
\]

\paragraph{(ii) Bhattacharyya Coefficient (BC) and its
Complement}\label{ii-bhattacharyya-coefficient-bc-and-its-complement}

The \textbf{Bhattacharyya Coefficient} provides a geometric measure of
similarity:

\[
\text{BC}(p, p^{\mathrm{rev}}) = \sum_{i=1}^{K} \sqrt{ p_i \cdot p_{K+1 - i} }.
\]

Its complement, \(1 - \text{BC}\), can be interpreted as a
divergence-like measure that reflects geometric differences between the
original and reversed distributions. While not a metric, it is bounded
between 0 and 1, with 0 indicating perfect symmetry.

\paragraph{(iii) Overlap Coefficient
(OVL)}\label{iii-overlap-coefficient-ovl}

The \textbf{Overlap Coefficient}, sometimes referred to as the
Szymkiewicz--Simpson coefficient, captures the shared probability mass:

\[
\text{OVL}(p, p^{\mathrm{rev}}) = \sum_{i=1}^{K} \min\{ p_i, p_{K+1 - i} \}.
\]

The complement \(1 - \text{OVL}\) thus reflects the fraction of mass
that is not shared between the two distributions.

These measures differ in their interpretive emphasis: - TV focuses on
the total discrepancy. - BC captures geometric congruence. - OVL
captures the total shared area between the distributions.

In our simulation study (Section ), we compute all three diagnostics to
illustrate how marginal asymmetry correlates with asymmetry in
achievable Pearson correlation bounds. These overlap-based measures are
implemented in the R package \texttt{overlap} (Ridout \& Linkie, 2009),
which facilitates reproducible estimation and visualization.

Empirically, we find that all three metrics provide consistent,
interpretable diagnostics of distributional asymmetry. Using them
jointly provides a robust triangulation of symmetry breaking, especially
in applied contexts where ordinal variables are discretized from latent
traits or ratings.

\section{Implications, Issues, and
Questions}\label{implications-issues-and-questions}

\begin{itemize}
\tightlist
\item
\end{itemize}

\subsection{Interpretation of `Effect
Size'}\label{interpretation-of-effect-size}

\begin{itemize}
\tightlist
\item
  Revisiting Cohen's benchmarks in light of constrained bounds
  \(\leadsto\) Interpretation of `effect size'
  (\textcite{cohe02-applied})
\item
  Risk of over- or under-stating association.
\end{itemize}

\subsection{Comparison Across Studies /
Measures}\label{comparison-across-studies-measures}

\begin{itemize}
\tightlist
\item
  Non-comparable bounds across datasets or variable codings.
\item
  Need for bound-aware effect size metrics.
\end{itemize}

\subsection{Hypothesis Testing}\label{hypothesis-testing}

\begin{itemize}
\tightlist
\item
  t-test can misslead: Example from BES: (i) t-test sig; perm test not;
  (ii) visa versa
\item
  Why p-values based on {[}--1, 1{]} bounds may mislead \(\leadsto\)
  need for Adjusted null distributions
\item
  \textbf{Permutation distribution:} Generate full vectors from
  marginals, independently permute \(X\) and \(Y\), compute \(r\),
  repeat. Produces null distribution centered at 0.
  \emph{(Proofs-OrdinalVariables.md)}
\end{itemize}

\subsection{Rescaling}\label{rescaling}

So, what to do? \emph{One} idea: \textbf{linear resclaing} (preserving
zero)

\begin{itemize}
\tightlist
\item
  Adjusted null distributions (via permutation or simulation).
\end{itemize}

\subsection{Variance ``Explained''}\label{variance-explained}

\begin{itemize}
\tightlist
\item
  Interpretation of \(r^2\) is ambiguous for at least two reasons:
\end{itemize}

\begin{enumerate}
\def\labelenumi{\arabic{enumi}.}
\tightlist
\item
  ordcats are invariant to monotone transformation (variance is a
  \emph{canrdinal} property, simply not approriate for ordinal
  variables)
\item
  is \(r^2_{max}\) the \ldots. maximal `amount' of variance between the
  two ordcats? How is this quantity to be interpreted? How is
  \(r^2_{min}\) to be interpreted?
\end{enumerate}

\section{Monte Carlo Simulations}\label{sec-MCsim}

To investigate factors that affect the bounds of attainable values of r
and thus the strength of constraints on r exerted by marginal
distributions we use two different sets of data. Our first dataset is
simulation generated and covers in principle all possible combinations
of marginal distributions for two orcats with given numbers of
categories and given n.~We simulate orcats containing, respectively, 4,
5, 6, 7, 10 and 11 categories, these being by far the most common in
empirical data deriving from survey studies. For each of these we
generate a large number of variables, each with its own simulated
univariate distribution, and thus also its own form (modality, symmetry,
etc.), mean, variance, skew and kurtosis as follows. {[}insert
description, detailing n, how many univariate distributions we generate
for each number of categories; how to populate the categories, whether
we impose distributional characteristics such as normality, modality,
skew etc. Btw: my preference would be to impose as little as possible,
and derive characteristics such as mean, variance, skew, kurtosis from
the distributions generated. {]} We subsequently randomly select pairs
of simulated variables and determine for each rmin, rmax, C1 and C2.
From the resulting simulated pairs of orcats we delete the cases in
which one, or both of the variables have zero variance. This results in
a dataset with n= xxx {[}insert very large number{]} that forms the
basis of subsequent analyses. Further details are specified in Appendix
3. This simulated dataset constitutes, within the bounds of granularity
of the procedure, a random sample all possible combinations of
univariate distributions of pairs of orcats, and can therefore be used
as an adequate basis for investigating functional relationships between
sundry properties of pairs of orcats and the range of attainable values
of r for those pairs. Although this simulated dataset randomly
represents all possible combinations of univariate distributions of
pairs of orcats, it cannot be considered to be a random sample of pairs
of orcats that applied researchers would encounter in actual, empirical
data such as derived from, e.g., mass surveys. This is because, when
designing survey questions, researchers tend to avoid formulations that
they expect to yield little variance in the responses. Therefore, we use
a second dataset to investigate constraints on the attainable range of
values of r, that is derived from the British Election Study 2019.
{[}fn: insert details about BES2019{]}. In this study we identified xxx
orcats, which define 7503 pairs, for each of which we calculated
pair-wise distributional characteristics including rmin, rmax, C1 and
C2. Further details on this dataset are provided in Appendix 4.
Obviously, the orcats in BES2019 are not a random sample from the
contents of all socio-political mass surveys in Britain in 2019, let
alone of those in wider areas such as the developed western world, or in
other years. But the BES can sensibly be regarded as typical for such
studies, which include (national) election studies, general social
surveys, and many public opinion and marketing surveys. As such,
findings based on this dataset derived from the BES can sensibly be
regarded as typical for what applied researchers in other countries will
encounter in their survey data. These two datasets will be used to
address, in the subsequent sections, questions with regard to: -
Distributions of rmin, rmax and measures of overall constraint of r by
marginals - Factors affecting rmin, rmax and overall constraint of r -
Consequences of constrained correlations in applied research, including
for descriptive purposes and hypothesis testing - Possibilities for
reducing undesired consequences of constrained correlations

VH: \((r_{min}, r_{max})\)

\subsection{Effects of Marginal
Shapes}\label{effects-of-marginal-shapes}

\begin{itemize}
\tightlist
\item
  Shape-based elasticity of correlation bounds\\
\item
  Skewness and modality heavily influence the attainable range of \(r\).
\item
  Marginals with extreme concentration (e.g., most responses in one or
  two categories) sharply restrict the possible correlation range.
\end{itemize}

\subsection{Symmetry Breaking in Correlation
Bounds}\label{symmetry-breaking-in-correlation-bounds}

\begin{itemize}
\item
  When and why \(r_{\min} \neq -r_{\max}\)
\item
  Empirical indicators of bound asymmetry
\item
  When given marginal distributions are sufficiently asymmetric, the min
  and max bounds of \(r\) are not symmetric about zero.
\item
  The shape of the null distribution of \(r\) can be strongly skewed,
  especially in cases with unequal marginals.
\end{itemize}

When we talk about ``asymmetry'' in this context, we're referring to how
different the distribution of \(Y\) is from the distribution of \(X\).
Since \(X\) is kept uniform, any deviation of \(Y\) from uniformity
creates asymmetry between the distributions.

The \emph{total variation distance} (TV) quantifies this asymmetry:

\begin{itemize}
\tightlist
\item
  It measures how different Y is from being symmetric around its
  midpoint
\item
  When TV = 0: \(Y\) is symmetric (like \(X\))
\item
  When TV is large: \(Y\) is highly asymmetric compared to \(X\)
\end{itemize}

\section{Illustration with `Typical' Data:
BES}\label{illustration-with-typical-data-bes}

We do not claim this data is representative of data-sets of ordinal
variables used by social scientists. Rather, we view it as
\emph{typical} in that it is a real-world example possessing many of the
features that would raise issues\ldots{}

VH: \((r_{min}, r_{max})\) And compare with MC simulated data

\section{Discussion}\label{discussion}

\subsection{Comparing `Typical' Data (BES) with Simulated
Data}\label{comparing-typical-data-bes-with-simulated-data}

\subsection{Relation to Literature}\label{relation-to-literature}

!fullcite camb76-inequalities

Full citation:

\textcite{camb76-inequalities}

Grether, D. M. (1976). On the Use of Ordinal Data in Correlation
Analysis. American Sociological Review, 41(5), 908--912.
https://doi.org/10.2307/2094741

Labovitz, S. (1967). Some observations on measurement and statistics.
Social Forces, 46(2), 151--160. https://doi.org/10.2307/2574595

Labovitz, S. (1970). The assignment of numbers to rank order categories.
American Sociological Review, 35(3), 515--524.
https://doi.org/10.2307/2092993

Moss, J., \& Grønneberg, S. (2023). Partial Identification of Latent
Correlations with Ordinal Data. Psychometrika, 88(1), 241--252.
https://doi.org/10.1007/s11336-022-09898-y

O'Brien, R. M. (1979). The use of Pearson's r with ordinal data.
American Sociological Review, 44(5), 851--857.
https://doi.org/10.2307/2094532

Robitzsch, A. (2022). On the Bias in Confirmatory Factor Analysis When
Treating Discrete Variables as Ordinal Instead of Continuous. Axioms,
11(4), Article 4. https://doi.org/10.3390/axioms11040162

\textcite{frec51-tableaux}

\textcite{whit76-bivariate}

\textcite{hoef40-masstabinvariante}

\textcite{lehm66-concepts}

\textcite{tche80-inequalities}

\textcite{camb76-inequalities}

\textcite{shih92-evaluating}

\textcite{alwi91-reliability}; \textcite{brew05-values};
\textcite{flor04-empirical}; \textcite{gait80-measurement};
\textcite{gree97-effect}; \textcite{jaco99-levels};
\textcite{kend38-new}; \textcite{knok02-statistics};
\textcite{kreu07-survey}; \textcite{kros97-designing};
\textcite{labo70-assignment}; \textcite{olss79-maximum};
\textcite{spea04-proof}

\subsection{Recommendations}\label{recommendations}

\begin{itemize}
\tightlist
\item
  Bound-aware effect size interpretation.
\item
  Bound-aware hypothesis testing.
\item
  Possible inclusion in statistical software.
\end{itemize}

\section{Conclusion}\label{conclusion}

Summarize:

\begin{itemize}
\item
  Main theoretical findings.
\item
  Key applied consequences.
\item
  Summary: Theoretical bounds, symmetry-breaking conditions, and
  implications for applied research
\item
  An R package accompanies this paper, providing tools to compute
  max/min bounds and perform randomization-based hypothesis tests for
  fixed marginals. Details are available in the appendix and online.
\end{itemize}

Impacts standard methods:

\begin{itemize}
\tightlist
\item
  interpretation (e.g.~`effect size') and comparison
\item
  hyothesis testing
\end{itemize}

Recommendations for interpreting Pearson's \(r\) in applied contexts

\begin{itemize}
\tightlist
\item
  Reporting observed \(r\) in context of its attainable bounds.
\item
  Using our tools to test whether an observed \(r\) is unusually strong
  or weak given marginals.
\end{itemize}

\subsection{Extensions and Open
Questions}\label{extensions-and-open-questions}

\begin{itemize}
\tightlist
\item
  Role of entropy and joint uncertainty in bounding behavior\\
\item
  Possibility of adjusting or normalizing \(r\)
\item
  Relation to Copula Theory (Optional)

  \begin{itemize}
  \tightlist
  \item
    ``Although our setting is discrete, the problem is conceptually
    linked to copula-based modeling of dependence with fixed marginals.
    The connection is discussed as a direction for future work.''
  \end{itemize}
\item
  `Correction' strategies \(\leadsto\) \emph{rescaling}
\item
  Generalizations to higher dimensions \(\leadsto\) implications for
  correlation matricies and Factor Analysis, SEMs, etc.

  \begin{itemize}
  \tightlist
  \item
    \textbf{PCA and Factor Analysis}: Attenuated correlations may lead
    to \(\leadsto\) ractical risks: over-factoring (more factors
    estimated than `should be'), biased structural coefficients(?)

    \begin{itemize}
    \tightlist
    \item
      \textbf{SEM}: Biased estimation of structural paths if
      ordinal-level attenuation is ignored.
    \end{itemize}
  \end{itemize}
\end{itemize}

\subsubsection{Correlation Matricies}\label{correlation-matricies}

The consequences of using Pearson correlations with ordinal data extend
beyond descriptive association. In factor analysis and SEM, attenuated
correlations can yield distorted loadings, inflated factor numbers, and
biased structural coefficients. Empirical studies confirm this risk:
\citet{flor04-emperical} show that CFA models fitted with Pearson's
\(r\) among ordinal indicators produce biased estimates compared to
models using polychoric correlations; \citet{gree97-efffect} demonstrate
how the number of response categories alters apparent model fit. More
broadly, psychometric research emphasizes that ignoring the ordinal
nature of indicators can misrepresent latent structures
(\autocite{lord68-statis}; \autocite{muth84-generl};
\autocite{bart11-laten}). For survey-based political science and
sociology, this means that correlations reported between Likert items
are not only bounded by their marginals, but may also systematically
bias downstream modeling.

Since the possible values \(r\) may take on depend on the marginal
distributions of the two variables being correlated, and since this
reference interval may -- and often is -- quite different for different
pairs of variables, interpreting and working with correlation matricies
(alternatively variance-covariance marticies) must be paid speacial
attention. Such matricies form the basis for a number of statistical
techniques, from regression to Factor Analysis or Structual Equation
Modeling (CITES). In this section we present preliminary results
exploring the implications of our findings above for PCA, FA, SEM.

First, some preliminaries. It is well known that any correlation matrix
(using completely cases -- i.e.~when each pair of variables is observed
for the same set of respondents), \(S\), is positive semi-definite and
symmetric. Hence, \(S\) is invertable. However, \ldots{}

\textbf{Why PSD Matters for PCA:}

\begin{itemize}
\tightlist
\item
  \textbf{Mathematically}: Eigen-decomposition works for any symmetric
  \(\mathbf{S}\), PSD or not.
\item
  \textbf{Interpretively}: Negative eigenvalues break the ``variance
  explained'' story (variance cannot be negative).
\item
  \textbf{Software behavior}: Most PCA routines check PD/PSD; failure
  can cause warnings or dropped components. \emph{(PCA and FA PSD
  requirement.md)}
\end{itemize}

\textbf{If Input is Symmetric but not PSD:}

Example:\\
\[
\mathbf{A} = \begin{bmatrix} 1 & 2 \\ 2 & 1 \end{bmatrix},\quad \lambda = (3,-1),
\] is indefinite (one negative eigenvalue).

\begin{itemize}
\tightlist
\item
  PCA math still runs; eigenvalues/vectors computed.
\item
  Interpretation fails: \(\sqrt{\lambda_j}\) imaginary for
  \(\lambda_j<0\).
\item
  Some software refuses input; others run but return nonsensical
  variance metrics. \emph{(PCA and FA PSD requirement.md)}
\end{itemize}

\textbf{PSD/PD Needs in FA:}

\begin{itemize}
\tightlist
\item
  \textbf{ML FA}: Requires \(\mathbf{S}\) be PSD;
  \(\log \det \mathbf{S}\) undefined if indefinite.
\item
  \textbf{Principal-axis / Image FA}: Also break with negative
  eigenvalues (square roots fail). \emph{(PCA and FA PSD
  requirement.md)}
\end{itemize}

\textbf{SEM and PSD Requirements}

\paragraph{SEM Machinery}\label{sem-machinery}

Observed covariance \(\mathbf{S}\); implied covariance
\(\boldsymbol\Sigma(\boldsymbol\theta)\) from structural and measurement
models. \emph{(PCA and FA PSD requirement.md)}

\begin{center}\rule{0.5\linewidth}{0.5pt}\end{center}

\paragraph{Fit Functions and PD Needs}\label{fit-functions-and-pd-needs}

\begin{itemize}
\tightlist
\item
  \textbf{ML / GLS}: Require both \(\mathbf{S}\) and
  \(\boldsymbol\Sigma(\theta)\) PD.
\item
  \textbf{WLS/DWLS}: Need PD indirectly (weight matrix invertible).
\item
  \textbf{ULS}: Can start from indefinite \(\mathbf{S}\), but test
  statistics again require PD. \emph{(PCA and FA PSD requirement.md)}
\end{itemize}

\begin{center}\rule{0.5\linewidth}{0.5pt}\end{center}

\paragraph{What Happens if Not PSD}\label{what-happens-if-not-psd}

\begin{itemize}
\tightlist
\item
  Software errors out (``matrix is not positive definite'').
\item
  Log-determinant undefined; inversion fails.
\item
  Even if bypassed, optimisation steps hit complex/NaN values.
  \emph{(PCA and FA PSD requirement.md)}
\end{itemize}

\begin{center}\rule{0.5\linewidth}{0.5pt}\end{center}

\paragraph{Remedies in SEM}\label{remedies-in-sem}

\begin{itemize}
\tightlist
\item
  Nearest PD smoothing (Higham 2002).
\item
  Ridge adjustment.
\item
  Analyse raw data.
\item
  Check rounding issues. \emph{(PCA and FA PSD requirement.md)}
\end{itemize}

\begin{center}\rule{0.5\linewidth}{0.5pt}\end{center}

\textbf{Core Lessons Across PCA, FA, SEM}

\begin{enumerate}
\def\labelenumi{\arabic{enumi}.}
\tightlist
\item
  \textbf{Theoretical PSD}: Covariance matrices from real data are PSD
  by definition.
\item
  \textbf{Numerical Issues}: Sampling noise, rounding, or listwise
  deletion can yield small negative eigenvalues.
\item
  \textbf{PCA}: Can compute from indefinite matrix but interpretation
  breaks.
\item
  \textbf{FA and SEM}: Require PD input; will not run on indefinite
  matrices.
\item
  \textbf{Repairs}: Adjust to nearest PSD, ridge, or recompute from raw
  data.
\end{enumerate}

So, we consider linearly re-scaling the estimated \(r\) for each pair of
variables in the correlation matirx \(S\). Call this new rescaled
correlation matrix \(S'\). We seek to understand what the process of
pair-wise linear rescaling has on the PSD-ness and invertability of
\(S'\). This is importnant because without an invertable corrleation
matrix, several analysese scholars perform on \(S\) becuase questionable
(at best). Further, some estimators used in e.g.~FA require PSD (namely,
MLE).

The general result is that linear rescaling does not `break'
invertability (nor PSD) but does make invertability `worse' vizaviz the
condition number. {[}Discuss condition number here{]}

\newpage

\section{References}\label{references}

\printbibliography[heading=none]

\newpage

\appendix

\section{Appendix A: Mathematical
Proofs}\label{appendix-a-mathematical-proofs}

\subsection{Proof A.1}\label{proof-max-product}

\textbf{Proof of \hyperref[prop-max-product]{Proposition
\textcite{prop-max-product}}.}

\subsection{Proof A.2}\label{proof-cor-min}

\textbf{Proof of \hyperref[cor-min]{Corollary 1}.}

\section{Appendix B: Software Tools and Code
Listings}\label{appendix-b-software-tools-and-code-listings}

\begin{itemize}
\tightlist
\item
  Documentation and examples for the R package\\
\item
  Description of randomization testing features
\end{itemize}

\section{Appendix C: Additional Tables and
Simulations}\label{appendix-c-additional-tables-and-simulations}

\begin{itemize}
\tightlist
\item
  Supplementary empirical examples, simulation results
\end{itemize}

\section{Appendix D: Old Outline}\label{appendix-d-old-outline}

\begin{enumerate}
\def\labelenumi{\arabic{enumi}.}
\tightlist
\item
  \textbf{Introduction and Conceptual Framework}
\end{enumerate}

\begin{itemize}
\tightlist
\item
  Research motivation and objectives
\item
  Overview of limitations in current correlation testing with ordinal
\item
  Why does this matter? \(\leadsto\) consequences if these changing
  bounds are ignored
\end{itemize}

\begin{enumerate}
\def\labelenumi{\arabic{enumi}.}
\setcounter{enumi}{1}
\tightlist
\item
  \textbf{Theoretical Bounds on Correlation Coefficients}
\end{enumerate}

\begin{itemize}
\tightlist
\item
  Mathematical explanation of R-min and R-max
\item
  Elasticity and asymmetry due to marginal distributions

  \begin{itemize}
  \tightlist
  \item
    Atainable and non-atainable bounds
  \end{itemize}
\item
  Implications for standard significance tests
\item
  Note on Freshet bounds anomaly (appendix or highlighted sidebar)
\end{itemize}

\begin{enumerate}
\def\labelenumi{\arabic{enumi}.}
\setcounter{enumi}{2}
\tightlist
\item
  \textbf{Measurement Issues and Metrics of Constraint} {[}TITLE?{]}
\end{enumerate}

\begin{itemize}
\tightlist
\item
  Problems with interval assumptions in ordinal categorical data
\item
  Alternatives to traditional scaling (psychometric, text-based)
\item
  Concept of boundedness and its distortion effects
\item
  Proposed \emph{metrics of constraint}: Asymmetry; Aitchinson distance;
  Entropy(?); Total Variation(?)
\end{itemize}

\begin{enumerate}
\def\labelenumi{\arabic{enumi}.}
\setcounter{enumi}{3}
\tightlist
\item
  \textbf{Simulations and Monte Carlo Approaches}
\end{enumerate}

\begin{itemize}
\tightlist
\item
  Design of permutation and Monte Carlo simulations
\item
  Bias and error analysis under null hypotheses
\item
  Role of mean, skewness, variance in predicting constraint metrics
\end{itemize}

\begin{enumerate}
\def\labelenumi{\arabic{enumi}.}
\setcounter{enumi}{4}
\tightlist
\item
  \textbf{Empirical Illustration Using BES Data}
\end{enumerate}

\begin{itemize}
\tightlist
\item
  Justification for selection: a typical but not representative case
\item
  Application of shrinkage and constraint metrics to BES
\item
  Comparison to simulation outcomes
\end{itemize}

\begin{enumerate}
\def\labelenumi{\arabic{enumi}.}
\setcounter{enumi}{5}
\tightlist
\item
  \textbf{Implications for Applied Research and Tools}
\end{enumerate}

\begin{itemize}
\tightlist
\item
  Critique of statistical purism in significance testing

  \begin{itemize}
  \tightlist
  \item
    Critique of over-reliance on significance without consideration of
    boundedness or constraint metrics.
  \end{itemize}
\item
  Reframing effect size: from absolute value to bounded context

  \begin{itemize}
  \tightlist
  \item
    \(r^2\) interpretation under bounded conditions
  \end{itemize}
\item
  RE-SCALING:

  \begin{enumerate}
  \def\labelenumi{\arabic{enumi}.}
  \tightlist
  \item
    Map \(r_{min}\) to -1, etc.
  \item
    Map \(\hat{r}^2\) to \([0,1]\), the re-scale \(\hat{r}_{adj}\) by
    preserving sign(?)
  \end{enumerate}

  \begin{itemize}
  \tightlist
  \item
    Effect on PSD and invertability of adjusted correlation matrix
    (NOTE: only use complete cases here so as to not `cheat' --
    \textcite{jaec09-positive}; \textcite{mish09-note};
    \textcite{shan23-impact})
  \end{itemize}
\item
  Practical guidance for analysts: how to assess when elastic bounds
  affect inference.
\item
  Overview of accompanying software tools or R functions developed.
\end{itemize}

\begin{enumerate}
\def\labelenumi{\arabic{enumi}.}
\setcounter{enumi}{6}
\item
  \textbf{Open Questions}
\item
  \textbf{Conclusions}
\end{enumerate}





\end{document}
